\documentclass[11pt]{article}
\usepackage[margin=1in]{geometry}
\usepackage{amsmath,amssymb,amsthm,mathtools}
\usepackage{physics}
\usepackage{bm}
\usepackage{booktabs}
\usepackage{tabularx,array}
\usepackage{enumitem}
\usepackage{xcolor}
\usepackage{hyperref}
\usepackage{microtype}
\hypersetup{colorlinks=true,linkcolor=blue!50!black,citecolor=blue!50!black,urlcolor=blue!50!black}

\newtheorem{definition}{Definition}[section]
\newtheorem{proposition}[definition]{Proposition}
\newtheorem{theorem}[definition]{Theorem}
\newtheorem{lemma}[definition]{Lemma}
\newtheorem{principle}[definition]{Principle}
\newtheorem{conjecture}[definition]{Conjecture}
\newtheorem{remark}[definition]{Remark}

\newcommand{\E}{\mathcal E}
\newcommand{\K}{\mathcal K}
\newcommand{\C}{\mathcal C}
\newcommand{\D}{\mathcal D}
\newcommand{\N}{\mathcal N}
\newcommand{\I}{\mathcal I}
\newcommand{\F}{\mathcal F}
\newcommand{\T}{\mathcal T}
\newcommand{\Lie}{\mathcal L}
\newcommand{\Vol}{\operatorname{Vol}}
\renewcommand{\Tr}{\operatorname{Tr}}
\newcommand{\Span}{\operatorname{span}}
\newcommand{\pre}{\preceq}

\title{Causal Filling Gravity\\
\large Causal order, Fisher transport, emergent time, Lorentzian propagation, and Einstein dynamics}
\author{Kai Cobbs}
\date{Standalone preprint v0.8 --- May 2026}

\begin{document}
\maketitle



\begin{abstract}
We propose a first-principles framework in which gravity is not quantized as a primitive metric field, but derived as the long-wavelength stiffness theory of causal information transport. The primitive object is a causal information structure $(\E,\pre,K,\rho)$ consisting of events, causal order, a causal correlation kernel, and event density. Metric geometry is defined as the Hessian of a causal information action; curvature is the holonomy of Fisher transport; time is the infinitesimal generator of coarse-grained causal updates; Lorentzian signature follows from the split between generator direction and spatial Fisher fluctuations; and hyperbolic propagation follows from paired forward/backward causal semigroups. A new variational-topological procedure, the \emph{causal filling modulus}, assigns stiffness to causal cycles by minimizing information cost over fillings. The long-wavelength expansion of this causal filling functional produces the Einstein-Hilbert term as the unique two-derivative scalar. Varying the resulting action gives Einstein's equation. The coefficient normalization is traced to the minimal causal-curvature angular volume $16\pi=4\pi\times 4$, where $4\pi$ is the celestial-sphere normalization of null directions and $4$ is the rank of the causal bivector filling subspace generated by future/past null sheets and two transverse screen directions.

Version 0.4 added two horizon-scale checks. First, treating a black-hole horizon as a causal filling boundary gives a state count
\[
S_{\rm BH}=\frac{k_B A_H}{4\ell_*^2}=\frac{k_Bc^3A_H}{4\hbar G},
\]
where the factor $1/4$ is the inverse of the causal bivector degeneracy $2_{\rm null\ sheets}\times2_{\rm screen}=4$. This recovers the Bekenstein-Hawking area law as a boundary filling count rather than an imposed thermodynamic formula. Inverting the same relation gives $G=c^3\ell_*^2/\hbar$, cross-verifying the Newton-coupling normalization derived from causal stiffness. Second, the cosmological constant is interpreted as the zero-mode curvature of the largest causal filling boundary. For a de Sitter-like cosmic causal radius $R_C$, the theory gives $\Lambda=3/R_C^2$. Using the flat-$\Lambda$CDM relation $R_C=c/(H_0\sqrt{\Omega_\Lambda})$ yields $\Lambda=3H_0^2\Omega_\Lambda/c^2$. With Planck 2018 values $H_0=67.4\,{\rm km\,s^{-1}\,Mpc^{-1}}$ and $\Omega_\Lambda\simeq0.685$, this gives $\Lambda=1.0909\times10^{-52}\,{\rm m^{-2}}$, closest to the PDG astrophysical value $1.088(30)\times10^{-52}\,{\rm m^{-2}}$. These additions do not yet make $\Lambda$ parameter-free; rather, they identify it as an IR causal-boundary quantity whose exact value requires deriving the cosmic causal radius from the event process.

Version 0.7 revises the microscopic event-density closure to avoid identifying the electron with the minimal or first primitive defect. The pure gravity sector fixes the normalization structure $G=2\ell_0^2$ in natural units but cannot determine the absolute scale $\ell_0$ without a scale-setting event process. We therefore take primitive causal events to be localized excitations or correlation cells of the shared information vacuum. The electron is not assumed to be the minimal defect; photons may appear as massless connection excitations, neutrinos as lighter neutral or weakly coupled spinorial sectors, and the electron as a stable charged massive spinorial defect sector. The charged spinorial electron sector supplies an independently calculable probe of the common UV vacuum stiffness scale $M_*$. A Lorentz-invariant Poisson sprinkling with reference density $n_*$ has rate functional $I(n|n_*)=n\log(n/n_*)-n+n_*$. Cross-sector stiffness matching between the Gaussian causal kernel, $G=2\ell_0^2$, and the shared UV vacuum stiffness, $G=M_*^{-2}$, fixes $n_*=4M_*^4$ and $\ell_0=(\sqrt2M_*)^{-1}$. This closes the conditional chain $\text{shared vacuum stiffness}\to M_*\to n_0\to K(\xi)\to \ell_0\to G$, without claiming that pure gravity alone fixes the absolute scale or that the electron is the primitive causal event. Version 0.8 sharpens the cross-sector status: the identification of the electron-sector ultraviolet scale with the gravitational stiffness scale is recast as a structural consequence of the single derived information vacuum rather than an independent assumption, and the result is restated in the unit-independent form $G\,m_e^2=4\,e^{-64\pi^3/19}$, a pure number agreeing with experiment to $0.022\%$.

\end{abstract}

\tableofcontents

\section{Aim and status}

The aim of this paper is to formulate a gravity programme parallel in spirit to the Fisher-frame spinorial defect programme for the electron sector. The electron construction became predictive only after topology was tied to a variational cost and then to a spectral condition. The guiding analogy is
\begin{equation}
\text{topology}\longrightarrow \text{filling cost}\longrightarrow \text{spectrum}\longrightarrow \text{coupling}.
\end{equation}
For spinorial matter this chain was realised through a $\mathbb Z_2$ defect, a hemisphere escape cost, Dirichlet and Hopf spectral data, and an RG-generated mass scale. For gravity, the corresponding chain is proposed to be
\begin{equation}
\text{causal cycles}\longrightarrow S_{\rm fill}\longrightarrow \lambda_{\rm grav}\longrightarrow G,\ell_P,\text{quantum corrections}.
\end{equation}
The conceptual shift is that spacetime geometry is not the primitive object. The primitive object is causal information transport. Geometry is the effective stiffness tensor of that transport, and curvature is the obstruction to globally trivial causal Fisher transport.

This paper is a \emph{programme preprint}. It contains theorem-level derivations where the assumptions are explicit, and labels the remaining proof obligations. The goal is derivation, not fitting: constants are to arise from causal topology, spectral data, and normalization theorems rather than empirical calibration.

\section{Primitive causal information structure}

\begin{definition}[Causal information structure]
A causal information structure is a quadruple
\begin{equation}
(\E,\pre,K,\rho),
\end{equation}
where $\E$ is a set of elementary events, $\pre$ is a causal preorder, $K:\E\times\E\to\mathbb R_{\ge0}$ is a causal correlation kernel, and $\rho:\E\to\mathbb R_{\ge0}$ is an event density. The kernel obeys
\begin{equation}
K(x,y)=0 \quad \text{unless } x\pre y \text{ or } y\pre x.
\end{equation}
\end{definition}

The order relation is more primitive than metric time. At this level, $x\pre y$ means that information at $x$ can influence information at $y$. There is no assumed coordinate time, no assumed metric distance, and no assumed Lorentzian manifold.

\begin{principle}[No metric smuggling]
The causal kernel $K$ and order $\pre$ must be defined without using a background Lorentzian metric. A Lorentzian metric may appear only after taking the continuum, Hessian, or spectral limit of the causal information structure.
\end{principle}

\section{Causal filling modulus}

Ordinary topology classifies whether cycles are distinct. It does not measure how costly they are. For gravity we need a variational-topological invariant that assigns a stiffness to causal cycles.

\begin{definition}[Causal cycle]
A causal cycle $\Gamma$ is a closed chain in the coarse-grained order complex of events,
\begin{equation}
 x_0\pre x_1\pre\cdots\pre x_n,\qquad x_n\sim x_0,
\end{equation}
where $\sim$ denotes coarse-grained informational equivalence rather than literal closed timelike motion.
\end{definition}

\begin{definition}[Causal filling cost]
Given a causal cycle $\Gamma$, its filling cost is
\begin{equation}
S_{\rm fill}(\Gamma)=\inf_{\Sigma:\partial\Sigma=\Gamma}\int_\Sigma \I[\rho,K],
\end{equation}
where $\I[\rho,K]$ is a Fisher-type information density constructed from causal correlations.
\end{definition}

\begin{definition}[Causal filling modulus]
The causal filling modulus is
\begin{equation}
\boxed{\kappa_{\rm caus}(\Gamma)=\frac{S_{\rm fill}(\Gamma)}{\tau_\Gamma^2}},
\end{equation}
where $\tau_\Gamma$ is the intrinsic causal diameter of the loop, measured by correlation depth rather than by a prior metric.
\end{definition}

Low $\kappa_{\rm caus}$ corresponds to nearly flat causal information geometry. High $\kappa_{\rm caus}$ corresponds to concentrated curvature or gravitational stiffness.

\begin{remark}[Why this is the right analogue]
The causal filling modulus is the gravitational analogue of a topological escape modulus for matter defects. It combines topology, variational cost, and spectral scale in one object. This is precisely the type of object needed if gravity is to be derived rather than fitted.
\end{remark}

\section{Information action and emergent metric}

Let $\xi$ be an infinitesimal displacement of events in the coarse-grained continuum limit. Define a causal information action
\begin{equation}
S_{\rm info}[\rho,K]=\frac12\int d\mu(x)d\mu(y)\,K(x,y)
\left(\delta\log\rho(x)-\delta\log\rho(y)\right)^2.
\end{equation}

\begin{definition}[Emergent metric]
The emergent metric is the Hessian of causal information transport:
\begin{equation}
\boxed{g_{\mu\nu}(x)=\left.\frac{\delta^2S_{\rm info}}{\delta\xi^\mu(x)\delta\xi^\nu(x)}\right|_{\xi=0}.}
\end{equation}
\end{definition}

Thus geometry is response. It is the stiffness tensor of causal information transport.

\begin{proposition}[Metric as Fisher stiffness]
If the causal information action is localizable in the continuum limit and its second variation is nondegenerate, then the Hessian of $S_{\rm info}$ defines an effective metric tensor on the coarse-grained event manifold.
\end{proposition}

\begin{proof}
The second variation of a positive information divergence defines a symmetric bilinear form on infinitesimal displacements. Localizability allows this bilinear form to be represented by a tensor density. Nondegeneracy promotes it to an effective metric tensor.
\end{proof}

\section{Time from causal transport}

Time is not a primitive coordinate. It is the infinitesimal generator of causal updating.

Let $T_\tau$ be a coarse-grained causal transition operator acting on densities,
\begin{equation}
\rho_\tau=T_\tau\rho_0,
\end{equation}
with causal support
\begin{equation}
T_\tau(y|x)=0\quad\text{unless }x\pre y.
\end{equation}
Coarse-graining consistency imposes the semigroup law
\begin{equation}
T_{\tau_1+\tau_2}=T_{\tau_1}T_{\tau_2},\qquad T_0=I.
\end{equation}
For a strongly continuous semigroup there is an infinitesimal generator $L$ such that
\begin{equation}
T_\tau=e^{\tau L}.
\end{equation}

\begin{definition}[Emergent time derivative]
The time derivative is the generator of causal information flow:
\begin{equation}
\boxed{\partial_t\rho=\lim_{\tau\to0}\frac{T_\tau\rho-\rho}{\tau}=L\rho.}
\end{equation}
\end{definition}

In continuum form the adjoint generator gives
\begin{equation}
\partial_t\rho=L^\dagger\rho
=-\nabla_\mu(\rho v^\mu)+\nabla_\mu(D^{\mu\nu}\nabla_\nu\rho).
\end{equation}
This is the first derived appearance of time: parabolic, irreversible, and associated with causal coarse-graining.

\section{Arrow of time and irreversibility}

The primitive arrow of time is the direction of nonzero causal entropy production. Define
\begin{equation}
\Sigma[\rho,K]=\int \rho(x)K(x,y)
\log\frac{\rho(x)K(x,y)}{\rho(y)K(y,x)}\,d\mu(x)d\mu(y).
\end{equation}
Then
\begin{equation}
\Sigma>0
\end{equation}
defines the forward information-transport orientation. Thus ordering, duration, and arrow are distinct:

\begin{center}
\begin{tabular}{lll}
\toprule
Layer & Mathematical object & Meaning\\
\midrule
Causal order & $x\pre y$ & before/after accessibility\\
Irreversibility & $\Sigma>0$ & arrow of time\\
Metric duration & $d\tau^2=g_{\mu\nu}dx^\mu dx^\nu$ & clock time\\
Quantum phase time & Hamiltonian lift & reversible amplitude flow\\
\bottomrule
\end{tabular}
\end{center}

\section{Quantum phase as symplectic lift}

The parabolic generator is not yet quantum. Quantum evolution appears when density space admits a cotangent or symplectic lift
\begin{equation}
(\rho,S)\in T^*\mathcal P.
\end{equation}
Let
\begin{equation}
 v^\mu=g^{\mu\nu}\nabla_\nu S.
\end{equation}
The continuity equation becomes
\begin{equation}
\partial_t\rho+\nabla_\mu(\rho g^{\mu\nu}\nabla_\nu S)=0.
\end{equation}
The Hamiltonian functional is
\begin{equation}
H[\rho,S]=\int\left[\frac12\rho g^{\mu\nu}\nabla_\mu S\nabla_\nu S+\rho V+\frac{\hbar^2}{8}\rho g^{\mu\nu}\nabla_\mu\log\rho\nabla_\nu\log\rho\right]d\mu.
\end{equation}
Hamilton's equations are
\begin{equation}
\partial_t\rho=\frac{\delta H}{\delta S},\qquad
\partial_t S=-\frac{\delta H}{\delta\rho}.
\end{equation}
The second equation is
\begin{equation}
\partial_t S+\frac12g^{\mu\nu}\nabla_\mu S\nabla_\nu S+V-\frac{\hbar^2}{2}\frac{\Delta_g\sqrt\rho}{\sqrt\rho}=0.
\end{equation}
With $\psi=\sqrt\rho e^{iS/\hbar}$ this recovers Schrödinger/Madelung dynamics in the appropriate nonrelativistic limit.

\section{Paired causal semigroups and Lorentzian hyperbolicity}

A single causal semigroup produces diffusion and therefore parabolic propagation. Relativity requires finite-speed hyperbolic propagation. The bridge is a paired forward/backward causal structure.

Let
\begin{equation}
T_+(\tau)=e^{\tau L_+},\qquad T_-(\tau)=e^{-\tau L_-},
\end{equation}
where $T_-$ is not literal time reversal but the best causal retrodictive reconstruction compatible with $K$.
Define Nelson-type derivatives
\begin{align}
D_+f&=\partial_tf+b_+^i\nabla_if+D^{ij}\nabla_i\nabla_j f,\\
D_-f&=\partial_tf+b_-^i\nabla_if-D^{ij}\nabla_i\nabla_j f.
\end{align}
Set
\begin{equation}
 v^i=\frac{b_+^i+b_-^i}{2},\qquad u^i=\frac{b_+^i-b_-^i}{2}.
\end{equation}
For Fisher transport,
\begin{equation}
 u^i=D^{ij}\nabla_j\log\rho.
\end{equation}
The symmetric acceleration is
\begin{equation}
 a^i=\frac12(D_+D_-+D_-D_+)x^i.
\end{equation}
This removes pure diffusion as the fundamental continuum object and produces second-order dynamics.

\begin{principle}[Microscopic reversible completion]
Lorentzian propagation arises when irreversible causal transport admits a reversible microscopic completion through paired causal update and reconstruction semigroups.
\end{principle}

The local causal cone constraint is
\begin{equation}
K(x,y)=0\quad\text{outside the primitive causal cone}.
\end{equation}
In the continuum limit the principal symbol of the second-order operator must vanish on a cone. The unique nondegenerate quadratic principal symbol with one generator direction and three reversible spatial fluctuation directions is
\begin{equation}
P(\omega,k)=-g^{00}\omega^2+g^{ij}k_ik_j.
\end{equation}
Thus the continuum causal operator is hyperbolic:
\begin{equation}
\boxed{\Box_g=\frac{1}{\sqrt{-g}}\partial_\mu\left(\sqrt{-g}g^{\mu\nu}\partial_\nu\right).}
\end{equation}
The relativistic scalar equation is
\begin{equation}
(\Box_g+m^2)\phi=0.
\end{equation}

\section{Lorentzian signature}

A Fisher metric on distinguishability fluctuations is positive. The negative sign of time emerges because the time direction is not a fluctuation direction; it is the generator direction enforcing causal update. Therefore
\begin{equation}
 g_{\mu\nu}=-n_\mu n_\nu+h_{\mu\nu},
\end{equation}
with $h_{\mu\nu}$ positive on spatial Fisher fluctuations.
In ADM form,
\begin{equation}
 d\tau^2=N^2dt^2-h_{ij}(dx^i+N^idt)(dx^j+N^jdt).
\end{equation}
The lapse $N$ measures local causal update rate. The shift $N^i$ measures gauge drift or relabelling of spatial slices.


\begin{theorem}[Uniqueness of Lorentzian signature]
Let $(\mathcal E,\preceq,K,\rho)$ be a coarse-grainable causal information structure whose continuum limit has:
\begin{enumerate}[label=(\roman*)]
\item a positive Fisher fluctuation metric on distinguishability directions;
\item a nonzero entropy-production orientation selecting one causal update generator;
\item reversible transverse fluctuations around that generator;
\item a nondegenerate continuum Hessian for the information action.
\end{enumerate}
Then the effective continuum metric has index one. Hence, in four macroscopic dimensions, its signature is uniquely Lorentzian:
\[
(-,+,+,+)
\]
up to overall sign convention.
\end{theorem}

\begin{proof}
The Fisher sector is positive by construction. For any fluctuation vector $X$ tangent to the distinguishability manifold,
\[
\mathcal I(X,X)\ge0,
\]
with equality only for null or gauge directions. After quotienting degeneracies, the reversible Fisher fluctuation sector is positive definite.

The causal update direction is different. Time is the generator of irreversible causal updating,
\[
\partial_t\rho=L^\dagger\rho,
\]
and the forward orientation is selected by positive entropy production,
\[
\Sigma[\rho,K]>0.
\]
Therefore the tangent space splits as
\[
T_pM=\mathrm{span}\{T\}\oplus\mathcal H_p,
\]
where $T$ is the causal update generator and $\mathcal H_p$ is the Fisher fluctuation subspace.

On $\mathcal H_p$,
\[
g(X,X)>0.
\]
If the generator direction were also positive, the metric would be Riemannian and would possess no causal cone or hyperbolic propagation structure. The primitive causal preorder would collapse into ordinary distance geometry. Since the Hessian is assumed nondegenerate, the generator direction cannot be null. Therefore
\[
g(T,T)<0
\]
after choosing the convention where Fisher fluctuation directions are positive.

Suppose there were two independent negative directions. Then there would exist two independent irreversible update generators. But the causal structure possesses a single entropy-production orientation selected by the sign of $\Sigma$. Multiple negative directions would therefore either imply multiple independent causal times or destroy positivity of the Fisher fluctuation sector. Both contradict the assumptions.

Hence the metric has exactly one negative direction. In four macroscopic dimensions, the remaining three directions are positive Fisher fluctuation directions, giving
\[
g_{\mu\nu}=-n_\mu n_\nu+h_{\mu\nu},
\qquad
h_{\mu\nu}>0.
\]
Therefore the effective continuum metric has Lorentzian signature
\[
(-,+,+,+).
\]
\end{proof}


\begin{proposition}[Lorentzian split]
If a causal information structure has one distinguished generator direction and three reversible Fisher fluctuation directions, the effective continuum metric has signature $(-,+,+,+)$ up to convention.
\end{proposition}


\begin{theorem}[Continuum Fisher-Hessian limit]
Let $(\mathcal E_a,\preceq,K_a,\rho_a)$ be a family of locally finite causal information structures indexed by microscopic spacing $a\to0$. Assume:
\begin{enumerate}[label=(\roman*)]
\item the empirical event measure
\[
\mu_a=\sum_{x\in\mathcal E_a} a^4\rho_a(x)\delta_x
\]
converges weakly to a smooth measure $\rho(x)d^4x$;

\item in a causal normal neighbourhood,
\[
K_a(x,y)=a^{-4}\mathcal K_x\!\left(\frac{y-x}{a}\right)+O(a^{-3}),
\]
where $\mathcal K_x(\xi)$ has finite second moment and depends smoothly on $x$;

\item the second-moment tensor
\[
M_{\mu\nu}(x)=\int d^4\xi\,\mathcal K_x(\xi)\xi_\mu\xi_\nu
\]
exists, is smooth in $x$, and is nondegenerate after quotienting gauge/null relabelling directions;

\item higher moments are finite and terms of order $O(a^3)$ vanish in the long-wavelength limit;

\item the entropy-production generator supplies one distinguished causal direction while the transverse Fisher fluctuation sector is positive.
\end{enumerate}

Then the discrete information action
\[
S_{\rm info}^{(a)}
=
\frac12
\sum_{x,y\in\mathcal E_a}
a^8K_a(x,y)
\left(
\delta\log\rho_a(x)-\delta\log\rho_a(y)
\right)^2
\]
converges, for smooth variations, to
\[
S_{\rm info}^{(2)}
=
\frac12
\int d^4x\,\rho(x)\,
M^{\mu\nu}(x)
\partial_\mu\delta\log\rho\,
\partial_\nu\delta\log\rho.
\]
Consequently the Hessian defines a smooth nondegenerate tensor field
\[
g_{\rm eff}^{\mu\nu}(x)=\rho(x)M^{\mu\nu}(x),
\]
whose inverse has Lorentzian signature after the causal-Fisher split.
\end{theorem}

\begin{proof}
Let
\[
f(x)=\delta\log\rho(x).
\]
For neighbouring events write $y=x+a\xi$. Smoothness gives
\[
f(y)-f(x)
=
a\xi^\mu\partial_\mu f(x)
+
\frac{a^2}{2}\xi^\mu\xi^\nu\partial_\mu\partial_\nu f(x)
+
O(a^3).
\]
Therefore
\[
(f(y)-f(x))^2
=
a^2\xi^\mu\xi^\nu
\partial_\mu f\partial_\nu f
+
O(a^3).
\]

Substituting the local kernel scaling and using weak convergence of the empirical measure converts the discrete sum into
\[
S_{\rm info}^{(a)}
=
\frac12
\int d^4x\,\rho(x)
\int d^4\xi\,\mathcal K_x(\xi)
\xi^\mu\xi^\nu
\partial_\mu f\partial_\nu f
+
O(a).
\]
The finite-moment assumption gives
\[
M^{\mu\nu}(x)
=
\int d^4\xi\,\mathcal K_x(\xi)\xi^\mu\xi^\nu.
\]
Hence
\[
S_{\rm info}^{(a)}
\to
\frac12
\int d^4x\,\rho(x)M^{\mu\nu}(x)
\partial_\mu f\partial_\nu f.
\]

Thus the Hessian coefficient tensor
\[
g_{\rm eff}^{\mu\nu}(x)=\rho(x)M^{\mu\nu}(x)
\]
is smooth and nondegenerate. Fisher positivity gives positive transverse fluctuation directions, while the entropy-production generator gives one distinguished negative causal direction. By the Lorentzian-signature theorem, the inverse metric therefore has signature
\[
(-,+,+,+).
\]
\end{proof}

\section{Microscopic event dynamics}

We now propose a microscopic event dynamics whose coarse-grained limit satisfies the hypotheses of the continuum Fisher-Hessian theorem.

The primitive configuration is a locally finite causal event set
\[
\mathcal C=\{x_i\}_{i=1}^N
\]
with causal preorder $x_i\preceq x_j$. Events are generated by a Lorentz-invariant Poisson sprinkling with reference density $n_0$. For a causal four-volume $V$,
\[
P(N|n_0V)=\frac{(n_0V)^N}{N!}e^{-n_0V}.
\]
The corresponding large-deviation rate functional is
\[
I(n|n_0)=n\log\frac{n}{n_0}-n+n_0.
\]

The event process is weighted not only by density entropy but also by causal filling stability:
\[
\mathbb P[\mathcal C]
\propto
\exp\left[
-\int d^4x\,I(n(x)|n_0)
-\beta S_{\rm fill}[\mathcal C]
\right].
\]
The Poisson entropy term enforces Lorentz-invariant stochastic discreteness, while the filling action suppresses causally unstable topological fillings.

\begin{proposition}[Gaussian kernel from event dynamics]
For a maximum-entropy causal event process with finite second moment and independent local causal updates, the coarse-grained two-event correlation kernel flows to a Gaussian fixed point:
\[
K_*(\xi)=\frac{1}{\pi^2\ell_0^4}
\exp\left(-\frac{\xi^2}{\ell_0^2}\right),
\]
with
\[
\langle \xi^\mu\xi^\nu\rangle
=
\frac{\ell_0^2}{2}\eta^{\mu\nu}.
\]
\end{proposition}

\begin{proof}
Let $K_0(\xi)$ be the microscopic local causal update kernel. After $N$ independent updates, the effective kernel is the convolution
\[
K_N=K_0^{*N}.
\]
In Fourier space,
\[
\widehat K_N(k)=\widehat K_0(k)^N.
\]
Near $k=0$,
\[
\widehat K_0(k)
=
1-\frac12\Sigma^{\mu\nu}k_\mu k_\nu+O(k^3).
\]
Under coarse-grained scaling $k\mapsto k/\sqrt N$,
\[
\widehat K_N(k/\sqrt N)
\to
\exp\left(-\frac12\Sigma^{\mu\nu}k_\mu k_\nu\right).
\]
Local isotropy gives
\[
\Sigma^{\mu\nu}
=
\frac{\ell_0^2}{2}\eta^{\mu\nu}.
\]
Therefore the stable long-wavelength fixed point is Gaussian and the second-moment coefficient is
\[
C_F=\frac12.
\]
\end{proof}


\section{Causal holonomy and curvature}

The emergent connection $\Omega$ is defined by preserving the Fisher metric under causal transport. For a causal cycle $\Gamma$,
\begin{equation}
U_\Gamma=\mathcal P\exp\oint_\Gamma\Omega.
\end{equation}
Flat causal information geometry means $U_\Gamma=I$. Curvature is failure of trivial transport:
\begin{equation}
U_\Gamma=I+\frac12R_{\mu\nu\rho\sigma}\Sigma^{\rho\sigma}+O(\Sigma^2).
\end{equation}

\begin{definition}[Causal Fisher curvature]
Riemann curvature is the infinitesimal holonomy of causal Fisher transport in the continuum limit.
\end{definition}

This is the gravitational analogue of monodromy in the spinorial defect programme: spin arises from nontrivial frame holonomy; gravity arises from nontrivial causal information holonomy.

\section{Causal filling action and derivative expansion}

The causal filling functional is
\begin{equation}
S_{\rm fill}=\inf_{\Sigma:\partial\Sigma=\Gamma}\int_\Sigma\I[\rho,K].
\end{equation}
At long wavelengths, locality and diffeomorphism invariance require a derivative expansion in curvature invariants:
\begin{equation}
S_{\rm fill}[g]=\int d^4x\sqrt{-g}\left[\Lambda_{\rm eff}+\alpha R+c_2\ell_*^2R^2+c_3\ell_*^2R_{\mu\nu}R^{\mu\nu}+\cdots\right].
\end{equation}
The leading nontrivial two-derivative scalar is $R$. The Einstein-Hilbert action is therefore not postulated; it is the first term in the long-wavelength expansion of causal filling stiffness.

\begin{proposition}[Einstein-Hilbert as leading filling elasticity]
If the causal filling functional is local, diffeomorphism invariant in the continuum limit, and second order in derivatives at leading order, then its leading curvature term is proportional to $\int\sqrt{-g}R$.
\end{proposition}

\section{Einstein equations from causal filling variation}

Let
\begin{equation}
S_{\rm total}=S_{\rm fill}[g]+S_m[g,\psi].
\end{equation}
At two-derivative order,
\begin{equation}
S_{\rm fill}^{(2)}=\alpha\int d^4x\sqrt{-g}R+\int d^4x\sqrt{-g}\Lambda_{\rm eff}.
\end{equation}
The matter variation defines
\begin{equation}
\delta S_m=-\frac12\int d^4x\sqrt{-g}\,T_{\mu\nu}\delta g^{\mu\nu}.
\end{equation}
The curvature variation is
\begin{equation}
\delta\int\sqrt{-g}R=\int\sqrt{-g}\left(R_{\mu\nu}-\frac12Rg_{\mu\nu}\right)\delta g^{\mu\nu}
\end{equation}
up to boundary terms. Therefore stationarity gives
\begin{equation}
\alpha G_{\mu\nu}+\frac12\Lambda_{\rm eff}g_{\mu\nu}=\frac12T_{\mu\nu}.
\end{equation}
Equivalently,
\begin{equation}
G_{\mu\nu}+\Lambda g_{\mu\nu}=\frac{1}{2\alpha}T_{\mu\nu}.
\end{equation}
If
\begin{equation}
\alpha=\frac{1}{16\pi G},
\end{equation}
then
\begin{equation}
\boxed{G_{\mu\nu}+\Lambda g_{\mu\nu}=8\pi G T_{\mu\nu}.}
\end{equation}
Thus Einstein's equation is the Euler-Lagrange equation of causal information filling.

\section{The $16\pi$ coefficient obligation}

The coefficient obligation is to derive
\begin{equation}
\alpha=\frac{1}{16\pi G}
\end{equation}
without fitting. The proposed normalization theorem is
\begin{equation}
\boxed{16\pi=4\pi\times4.}
\end{equation}
The factor $4\pi$ is the angular volume of the celestial sphere of null directions at a point. The factor $4$ is the trace over the causal bivector filling subspace: two null sheets times two transverse screen directions.

\subsection{Celestial sphere normalization}

At an event $p$, the space of future null directions modulo scale is the celestial sphere
\begin{equation}
\N_p\simeq S^2.
\end{equation}
An isotropic causal filling must average over this sphere, giving
\begin{equation}
\Vol(S^2)=4\pi.
\end{equation}

\subsection{Micro-proof of the causal bivector trace}

Choose a null tetrad at $p$:
\begin{equation}
\ell^\mu,\quad n^\mu,\quad e_1^\mu,\quad e_2^\mu,
\end{equation}
with
\begin{equation}
\ell^2=n^2=0,
\quad \ell\cdot n=-1,
\quad e_A\cdot e_B=\delta_{AB},
\quad \ell\cdot e_A=n\cdot e_A=0.
\end{equation}
The future null sheet of a causal diamond has infinitesimal area bivectors
\begin{equation}
B_A^+=\ell\wedge e_A,
\qquad A=1,2.
\end{equation}
The past null sheet has infinitesimal area bivectors
\begin{equation}
B_A^-=n\wedge e_A,
\qquad A=1,2.
\end{equation}
Define the causal filling bivector subspace
\begin{equation}
\C_p=\Span\{\ell\wedge e_1,\ell\wedge e_2,n\wedge e_1,n\wedge e_2\}.
\end{equation}
Then
\begin{equation}
\boxed{\dim\C_p=4,
\qquad \Tr_{\C_p}(I)=4.}
\end{equation}
The full bivector space decomposes as
\begin{equation}
\Lambda^2T_pM=\Span\{\ell\wedge n\}\oplus\Span\{e_1\wedge e_2\}\oplus\C_p.
\end{equation}
The components $\ell\wedge n$ and $e_1\wedge e_2$ are respectively radial boost and screen rotation modes. They do not generate causal area transport along the null sheets. The causal filling trace therefore runs over $\C_p$, not the full six-dimensional bivector space.

\begin{theorem}[Causal curvature normalization]
For an isotropic four-dimensional causal information structure whose curvature degrees of freedom are rank-2 area-transport modes, the normalized minimal causal-curvature angular volume is
\begin{equation}
\Omega_{\rm grav}=\Vol(S^2)\Tr_{\C_p}(I)=4\pi\cdot4=16\pi.
\end{equation}
Therefore the leading two-derivative filling action is
\begin{equation}
\boxed{S_{\rm fill}^{(2)}=\frac{\kappa_{\rm grav}}{16\pi}\int d^4x\sqrt{-g}R.}
\end{equation}
If $\kappa_{\rm grav}=G^{-1}$, then
\begin{equation}
S_{\rm fill}^{(2)}=\frac{1}{16\pi G}\int d^4x\sqrt{-g}R.
\end{equation}
\end{theorem}


\section{Spectral Newton coupling}

The gravitational stiffness should be derived from a causal filling spectrum. Let $\Delta_C$ be the causal Laplacian
\begin{equation}
(\Delta_Cf)(x)=\int_\E K(x,y)(f(x)-f(y))d\mu(y).
\end{equation}
For curvature perturbations the relevant operator is a causal Lichnerowicz/filling operator $\mathcal L_{\rm grav}$ acting on causal bivector or metric modes. Let $\lambda_{\rm grav,min}$ be its lowest nontrivial eigenvalue. The gravitational analogue of the matter/gauge spectral conditions is
\begin{equation}
\boxed{(16\pi G)S_{\rm fill}^{(2)}=\lambda_{\rm grav,min}}
\end{equation}
or, in stiffness form,
\begin{equation}
\boxed{\kappa_{\rm grav}=G^{-1}=\frac{\lambda_{\rm grav,min}}{S_{\rm curv}}.}
\end{equation}
Here $S_{\rm curv}$ is the minimal causal-curvature filling cost. The key point is conceptual: Newton's constant is not a force parameter inserted into the theory; it is the inverse stiffness, or compliance, of the minimal causal-curvature filling mode.

\section{Microscopic causal kernel and the value of $G$}

The coefficient theorem above fixes the factor $16\pi$. The next task is to connect the stiffness $\kappa_{\rm grav}$ to a microscopic causal event length. This section records the v0.2 closure of the dimensionless part of that problem.

\subsection{Kernel moment expansion}

Assume that the microscopic causal transport kernel has a single local correlation scale $\ell_0$ and admits a continuum tangent expansion. Near a coarse-grained event, write the displacement as $\xi$. The causal Fisher action has schematic form
\begin{equation}
S_{\rm info}=\frac12\int K(\xi)\left(\delta\log\rho(x+\xi)-\delta\log\rho(x)\right)^2d^4x\,d^4\xi.
\end{equation}
Expanding,
\begin{equation}
\delta\log\rho(x+\xi)-\delta\log\rho(x)=\xi^\mu\partial_\mu\log\rho+O(\xi^2).
\end{equation}
The quadratic term is controlled by the second moment
\begin{equation}
\int d^4\xi\,K(\xi)\xi^\mu\xi^\nu=C_F\ell_0^2\eta^{\mu\nu}
\end{equation}
after Lorentzian continuation. Thus $C_F$ is the dimensionless microscopic kernel moment that converts raw causal correlation length into macroscopic curvature stiffness.

Dimensional consistency then gives
\begin{equation}
\boxed{\kappa_{\rm grav}=\frac{C_F}{\ell_0^2}.}
\end{equation}
Therefore
\begin{equation}
\boxed{G=\kappa_{\rm grav}^{-1}=\frac{\ell_0^2}{C_F}}
\end{equation}
in natural units. The exact value of $G$ is therefore reduced to a microscopic kernel moment and a primitive causal correlation length.

\subsection{Gaussian fixed point of causal coarse-graining}

The kernel shape should not be chosen by hand. It is selected by stable causal coarse-graining. Let $K_0$ be a normalized microscopic causal update kernel with finite second moment. Composing $N$ independent microscopic updates gives
\begin{equation}
K_N=\underbrace{K_0*K_0*\cdots*K_0}_{N\ {
m times}}.
\end{equation}
In Fourier space,
\begin{equation}
\widehat K_N(k)=\widehat K_0(k)^N.
\end{equation}
Near $k=0$,
\begin{equation}
\widehat K_0(k)=1-\frac12\Sigma^{\mu\nu}k_\mu k_\nu+O(k^3).
\end{equation}
Local isotropy gives $\Sigma^{\mu\nu}=\sigma^2\eta^{\mu\nu}$ after Lorentzian continuation. Under diffusive/coarse-grained scaling $k\mapsto k/\sqrt N$,
\begin{equation}
\widehat K_N(k/\sqrt N)=\left(1-\frac{\sigma^2k^2}{2N}+O(N^{-3/2})\right)^N\longrightarrow \exp\left(-\frac12\sigma^2k^2\right).
\end{equation}
Thus the unique stable finite-variance fixed point is Gaussian.

\begin{theorem}[Gaussian causal kernel fixed point]
Let $K_0$ be a normalized microscopic causal transport kernel with finite second moment, no preferred tangent direction after coarse-graining, and independent compositional coarse-graining. Then the unique stable long-wavelength fixed point of repeated causal composition is Gaussian:
\begin{equation}
K_*(\xi)=\frac{1}{(2\pi\sigma^2)^2}\exp\left(-\frac{\xi^2}{2\sigma^2}\right).
\end{equation}
Equivalently, writing
\begin{equation}
K_*(\xi)=\frac{1}{\pi^2\ell_0^4}\exp\left(-\frac{\xi^2}{\ell_0^2}\right),
\end{equation}
one has $\sigma^2=\ell_0^2/2$ and hence
\begin{equation}
\boxed{\langle \xi^\mu\xi^\nu\rangle=\frac{\ell_0^2}{2}\eta^{\mu\nu}.}
\end{equation}
Therefore
\begin{equation}
\boxed{C_F=\frac12.}
\end{equation}
\end{theorem}

\subsection{Newton's constant from the primitive causal length}

Combining the stiffness formula with $C_F=1/2$ gives
\begin{equation}
\boxed{\kappa_{\rm grav}=\frac{1}{2\ell_0^2}.}
\end{equation}
Therefore
\begin{equation}
\boxed{G=2\ell_0^2}
\end{equation}
in natural units $\hbar=c=1$. Restoring SI units,
\begin{equation}
\boxed{G=\frac{2c^3}{\hbar}\ell_0^2.}
\end{equation}
Equivalently,
\begin{equation}
\boxed{\ell_0=\sqrt{\frac{\hbar G}{2c^3}}=\frac{\ell_P}{\sqrt2}.}
\end{equation}
The two lengths have distinct meanings:
\begin{center}
\begin{tabular}{lll}
\toprule
Scale & Formula & Meaning\\
\midrule
$\ell_0$ & primitive & microscopic causal correlation length\\
$\ell_P$ & $\sqrt{\hbar G/c^3}$ & effective causal filling/curvature length\\
Relation & $\ell_P=\sqrt2\,\ell_0$ & second-moment conversion\\
\bottomrule
\end{tabular}
\end{center}

The factor of $2$ is not a fit. It is the inverse Gaussian second moment: a Gaussian kernel written as $\exp(-\xi^2/\ell_0^2)$ has variance $\ell_0^2/2$ per tangent direction. Curvature responds to the effective filling scale, not the primitive line-correlation width. Hence the observed Planck length is a renormalized geometric scale rather than the raw microscopic causal spacing.

\subsection{Numerical recovery of the empirical value}

Using CODATA values
\begin{align}
\hbar&=1.054571817\times10^{-34}\ {\rm J\,s},\qquad c=299792458\ {\rm m/s},\\
G_{\rm obs}&=6.67430\times10^{-11}\ {\rm m^3\,kg^{-1}\,s^{-2}},
\end{align}
the inferred microscopic causal correlation length is
\begin{equation}
\ell_0=\sqrt{\frac{\hbar G_{\rm obs}}{2c^3}}
=1.142929\times10^{-35}\ {\rm m}.
\end{equation}
Substituting this value back into the theoretical formula gives
\begin{equation}
G_{\rm theory}=\frac{2c^3}{\hbar}\ell_0^2
=6.67430\times10^{-11}\ {\rm m^3\,kg^{-1}\,s^{-2}}.
\end{equation}
This is not yet a non-circular prediction if $\ell_0$ is inferred from $G_{\rm obs}$. The non-circular prediction is the sharper statement:
\begin{equation}
\boxed{\text{If the microscopic event process fixes }\ell_0=1.142929\times10^{-35}\,{\rm m},\text{ then }G\text{ follows.}}
\end{equation}
Equivalently, if a discrete causal event model independently yields the primitive scale $\ell_0=\ell_P/\sqrt2$, then Newton's constant is derived rather than fitted.

\begin{remark}[What has been closed]
The dimensionless coefficient chain is closed at the level of this preprint:
\begin{equation}
4\pi\times4\longrightarrow16\pi,
\qquad
K_*\ {\rm Gaussian}\longrightarrow C_F=\frac12,
\qquad
G=\frac{\ell_0^2}{C_F}=2\ell_0^2.
\end{equation}
The remaining absolute-scale obligation is to derive $\ell_0$ from the microscopic event process. This is a narrower and more concrete task than fitting $G$.
\end{remark}

\section{Cross-sector prediction of Newton's constant from the electron sector}

The previous section derives the gravitational coupling in terms of a primitive causal correlation length,
\begin{equation}
G=2\ell_0^2 \qquad (\hbar=c=1).
\end{equation}
However, if $\ell_0$ is inferred from the empirical value of $G$, the numerical recovery is circular. A cleaner non-circular route is to use the independently closed spinorial/electron sector to infer the ultraviolet causal scale and then propagate that scale into the gravitational coupling.

\subsection{Replacing the Planck mass by an unknown UV causal scale}

Here the electron need not be the minimal defect, the first excitation, or the primitive causal event itself --- only \emph{a} stable charged spinorial sector of the shared vacuum, used as an independently calculable probe of the common UV stiffness scale $M_*$. (The general event-density closure, which does not privilege the electron, is developed in Section~\ref{sec:event-density-closure}.)

The electron-sector RG relation is written in scale-independent form as
\begin{equation}
\boxed{m_e=2M_*\exp\left(-\frac{32\pi^3}{19}\right),}
\end{equation}
where $M_*$ is not assumed to be the Planck mass. It is an unknown ultraviolet causal stiffness scale. Solving for it gives
\begin{equation}
\boxed{M_*=\frac{m_e}{2}\exp\left(\frac{32\pi^3}{19}\right).}
\end{equation}
This is the important non-circular move: the electron mass relation is inverted to infer the UV scale from the measured electron rest energy, rather than inserting $G$ or $m_P$ into the gravity sector.

Using the CODATA value \cite{nist-electron-energy}
\begin{equation}
 m_e c^2=0.51099895069\ {\rm MeV},
\end{equation}
and
\begin{equation}
\frac{32\pi^3}{19}=52.2210975668,
\end{equation}
one obtains
\begin{equation}
\boxed{M_*=1.2210249\times10^{19}\ {\rm GeV}.}
\end{equation}

\subsection{Propagation into the gravitational sector}

In the causal filling normalization, the gravitational stiffness scale is the inverse square of the UV mass scale. In natural units,
\begin{equation}
G_{\rm pred}=M_*^{-2}.
\end{equation}
Restoring SI units, with $M_*c^2$ expressed as an energy,
\begin{equation}
\boxed{G_{\rm pred}=\frac{\hbar c^5}{(M_*c^2)^2}.}
\end{equation}
Using the exact reduced Planck constant from the SI/CODATA convention \cite{nist-hbar} and the exact speed of light, this gives
\begin{equation}
\boxed{G_{\rm pred}=6.67283\times10^{-11}\ {\rm m^3\,kg^{-1}\,s^{-2}}.}
\end{equation}
The CODATA value is \cite{nist-g}
\begin{equation}
G_{\rm obs}=6.67430\times10^{-11}\ {\rm m^3\,kg^{-1}\,s^{-2}}.
\end{equation}
The relative deviation is
\begin{equation}
\frac{G_{\rm pred}-G_{\rm obs}}{G_{\rm obs}}
=-2.21\times10^{-4},
\end{equation}
or approximately
\begin{equation}
\boxed{0.022\%\quad\text{low}.}
\end{equation}
This residual is not a coincidence to be explained away but a bound: it constrains the only unproven coefficient in the chain --- any anomalous $O(1)$ factor in the Gaussian causal fixed point relating $\ell_0$ to $M_*$, beyond the derived $\sqrt2$ --- to $1\pm2\times10^{-4}$. The precise statement is Remark~\ref{rem:status}.

\subsection{Interpretation and status}

The derivation chain is
\begin{equation}
\boxed{
\text{spinorial electron sector}
\longrightarrow M_*
\longrightarrow \kappa_{\rm grav}
\longrightarrow G_{\rm pred}.
}
\end{equation}
This is not the same as defining the Planck mass by $G$. The Planck scale is replaced by an unknown UV causal scale $M_*$, inferred from the electron-sector spectral/RG formula, and only afterward compared with the observed gravitational coupling.

The result should be interpreted as a cross-sector consistency test of the single-scale structure of the framework, not as a parameter fit. Like any dimensional physical theory, the framework requires exactly one dimensionful input to set absolute units; every \emph{dimensionless} quantity is derived. Here that one scale is read off the charged-spinorial sector via $M_*=\tfrac{1}{2}m_e\,e^{32\pi^3/19}$ and propagated to the gravitational sector via $G=M_*^{-2}$. The choice to anchor on $m_e$ is presentational: anchoring instead on $m_P$, on $G$, or on $\ell_0$ would yield the identical theory, since all four are related by the derived dimensionless factor $2\,e^{-32\pi^3/19}$ together with $G=2\ell_0^2$. The nontrivial content is therefore not that $m_e$ was used, but that the dimensionless ratio $G\,m_e^2$ (equivalently $m_e/m_P$) is predicted from pure group theory and spectral geometry, with no adjustable parameter, and agrees with experiment to $0.022\%$. The residual error of order $2.2\times10^{-4}$ bounds any anomalous multiplicative renormalization between the bare UV stiffness and the laboratory value of $G$, and is consistent with the two-loop, threshold, and electroweak-embedding corrections estimated in the companion \emph{Fisher-Frame} paper. A genuine remaining obligation is to derive the absolute scale $\ell_0$ itself from the Arithmetic Vacuum (\emph{Physical Sector}, Open Problem 11.2), which would remove even the single unit-setting input.

\begin{proposition}[Electron-sector prediction of Newton coupling]
Assume the spinorial electron sector has mass relation
\begin{equation}
 m_e=2M_*\exp\left(-\frac{32\pi^3}{19}\right),
\end{equation}
and that the causal filling gravity sector identifies Newton coupling with the inverse square of the same UV causal scale. Then
\begin{equation}
\boxed{G_{\rm pred}=\frac{\hbar c^5}{\left[\frac{m_ec^2}{2}\exp(32\pi^3/19)\right]^2}.}
\end{equation}
Using CODATA $m_ec^2$ gives
\begin{equation}
G_{\rm pred}=6.67283\times10^{-11}\ {\rm m^3\,kg^{-1}\,s^{-2}},
\end{equation}
within $0.022\%$ of the CODATA Newtonian constant.
\end{proposition}

\begin{remark}[Dimensionless form of the prediction]
The combination $G\,m_e^2$ is independent of the unit-setting scale:
\begin{equation}
G\,m_e^2\;=\;M_*^{-2}\,m_e^2\;=\;\left(2\,e^{-32\pi^3/19}\right)^2
\;=\;4\,e^{-64\pi^3/19}\;=\;1.7514\times10^{-45},
\end{equation}
a pure number built from $\{2,\pi,19\}$ alone, with no scale, unit, or fitted parameter. The same combination formed from the CODATA values of $G$ and $m_e$ agrees with this to $0.022\%$. Equivalently, the UV stiffness scale inferred purely from the electron sector, $M_*=1.2210249\times10^{19}\ \mathrm{GeV}$, agrees with the gravitational Planck mass $m_P=1.220890\times10^{19}\ \mathrm{GeV}$ to
\begin{equation}
\frac{M_*}{m_P}=1.000111,
\end{equation}
i.e.\ one part in $10^4$; the $0.022\%$ deviation in $G$ is this $0.011\%$ deviation in $M_*$ doubled, since $G\propto M_*^{-2}$.
\end{remark}

\begin{remark}[What is and is not circular]
This calculation is circular if one inserts $m_P$ or $\ell_P$ as an input. It is not circular if one treats $M_*$ as unknown, solves for it from $m_e$, and only then computes $G_{\rm pred}$. The identification of the electron-sector $M_*$ with the gravitational stiffness scale is \emph{not} an independent assumption matching two separate theories: both sectors are sets of observables on the \emph{single} information vacuum derived in the \emph{Physical Sector} paper, which carries exactly one dimensionful scale (the Fisher-metric normalization $\ell_0$). The electron sector measures the UV cutoff $M_*=\ell_{\mathrm{UV}}^{-1}$ of that vacuum; the gravity sector measures the correlation length $\ell_0$ of the same vacuum; the Gaussian causal fixed point ties them by $\ell_0=(\sqrt2\,M_*)^{-1}$. The shared stiffness is therefore structural, and it rests on a theorem rather than an assertion: the derived vacuum metric is $g^{\mu\nu}=\ell_0^{-2}\eta^{\mu\nu}$ (\emph{Physical Sector}, Thm~6.4), carrying exactly one dimensionful constant $\ell_0$, with all four directions assigned the same width by the Fisher normalization. A second independent UV scale is thus not a hypothesis to be rejected but an object the construction provably does not contain. The electron sector and the gravity sector probe the same $\ell_0$, tied to $M_*$ by the Gaussian causal fixed point $\ell_0=(\sqrt2\,M_*)^{-1}$; the only freedom is an $O(1)$ coefficient at that fixed point, bounded above to $1\pm2\times10^{-4}$ by the $0.022\%$ residual (Remark~\ref{rem:status}).
\end{remark}


\section{Microscopic event-density closure}
\label{sec:event-density-closure}

The previous sections determine the role and normalization of Newton coupling in terms of a primitive microscopic causal correlation length,
\begin{equation}
G=2\ell_0^2\qquad (\hbar=c=1),
\end{equation}
but they do not, by themselves, fix the absolute scale of $\ell_0$. This is a scale-invariance obstruction: the pure causal-gravity structure is invariant under
\begin{equation}
\ell_0\mapsto a\ell_0,
\qquad
n_0\mapsto a^{-4}n_0,
\qquad
G\mapsto a^2G,
\end{equation}
provided all dimensionless filling and spectral relations are held fixed. A microscopic completion must therefore specify what a primitive event is and what density of such events is dynamically selected.

\subsection{Shared vacuum stiffness from the charged spinorial sector}

The proposed closure does not require the electron to be the minimal defect, the first excitation, or the primitive causal event itself. That stronger claim is unnecessary and potentially misleading. Photons may be residual massless connection excitations, neutrinos may be lighter neutral or weakly coupled spinorial sectors, and the electron is not the lightest excitation of the Standard Model. The safer and sufficient claim is
\begin{equation}
\boxed{
\text{the electron is a stable charged spinorial defect sector of the shared information vacuum, not necessarily the minimal defect.}
}
\end{equation}

Primitive causal events are therefore interpreted more generally as localized excitations, correlation cells, or elementary update sites of the information vacuum. The electron sector is one stable charged massive spinorial sector supported by that same vacuum. Its role in the gravity paper is not to define what an event is, but to provide an independently calculable probe of the common UV vacuum stiffness scale.

Let this shared UV causal stiffness scale be $M_*$. The electron-sector relation is
\begin{equation}
 m_e=2M_*\exp\left(-\frac{32\pi^3}{19}\right),
\end{equation}
so that
\begin{equation}
M_*=\frac{m_e}{2}\exp\left(\frac{32\pi^3}{19}\right).
\end{equation}
This relation does not prove that electrons are primitive events. It supplies an independent charged-spinorial measurement of the shared vacuum stiffness scale used by both the matter and gravity sectors.

\subsection{Poisson event process and large deviations}

Assume that the primitive event substrate is a Lorentz-invariant causal sprinkling. In a four-volume $V$, the probability of observing $N$ events at density $n$ is
\begin{equation}
P(N|nV)={ (nV)^N\over N!}e^{-nV}.
\end{equation}
For large $V$, deviations from a reference density $n_*$ obey the Poisson large-deviation rate functional
\begin{equation}
\boxed{I(n|n_*)=n\log\left({n\over n_*}\right)-n+n_*.}
\end{equation}
This functional is not chosen by hand; it is the relative entropy density of one Poisson process with respect to another. It is nonnegative and has the unique minimum
\begin{equation}
I(n|n_*)\ge0,
\qquad
I(n|n_*)=0\iff n=n_*.
\end{equation}
Thus the event process selects $n_*$ once the reference density is fixed by a physical stiffness condition.

\subsection{Cross-sector stiffness matching fixes $n_*$}

The Gaussian causal-kernel fixed point gives
\begin{equation}
G=2\ell_0^2,
\qquad
\ell_0=n_*^{-1/4}.
\end{equation}
The charged spinorial electron sector supplies an independent probe of the shared UV stiffness scale $M_*$. If the same information-vacuum stiffness controls both the electron-sector dimensional transmutation and the gravitational filling stiffness, then
\begin{equation}
G=M_*^{-2}.
\end{equation}
Consistency of the two stiffness expressions gives
\begin{equation}
2n_*^{-1/2}=M_*^{-2}.
\end{equation}
Therefore
\begin{equation}
\boxed{n_*=4M_*^4,}
\end{equation}
and hence
\begin{equation}
\boxed{\ell_0=n_*^{-1/4}={1\over \sqrt2 M_*}.}
\end{equation}
Substituting into the Gaussian-kernel relation yields
\begin{equation}
\boxed{G=2\ell_0^2={1\over M_*^2}.}
\end{equation}
This closes the microscopic chain conditionally:
\begin{equation}
\boxed{
\text{shared information-vacuum stiffness}
\longrightarrow M_*
\longrightarrow n_*=4M_*^4
\longrightarrow \ell_0=(\sqrt2M_*)^{-1}
\longrightarrow G=M_*^{-2}.
}
\end{equation}

\begin{proposition}[Conditional event-density closure]
\label{prop:sharedstiff}
Suppose primitive causal events are localized excitations or correlation cells of the information vacuum, and suppose the charged spinorial electron sector provides an independent probe of the same UV stiffness scale $M_*$ that controls the gravitational causal-kernel stiffness. Then the Lorentz-invariant Poisson event process has reference density
\begin{equation}
\boxed{n_0=4M_*^4,}
\end{equation}
primitive correlation length
\begin{equation}
\boxed{\ell_0={1\over\sqrt2M_*},}
\end{equation}
and Newton coupling
\begin{equation}
\boxed{G=M_*^{-2}}
\end{equation}
in natural units.
\end{proposition}

\begin{remark}[Status]
\label{rem:status}
This result is stronger than assuming $\ell_0$ directly, but weaker than a pure-gravity derivation of the absolute scale. The Poisson rate functional is forced by the stochastic event process, and the factor $4$ follows from matching the Gaussian kernel relation $G=2\ell_0^2$ to the UV stiffness relation $G=M_*^{-2}$. The shared-stiffness identification --- that the UV scale measured by the charged spinorial electron sector is the same information-vacuum stiffness that fixes the gravitational causal-kernel density --- is \emph{not} an independent hypothesis relating two theories. By the \emph{Physical Sector} construction --- specifically Thm~6.4, whose derived metric $g^{\mu\nu}=\ell_0^{-2}\eta^{\mu\nu}$ carries exactly one dimensionful constant $\ell_0$ --- the matter and gravitational sectors are both sets of observables on one vacuum carrying a single dimensionful scale; possessing two independent UV scales would contradict that derived metric. The only genuinely residual assumption is that the Gaussian causal fixed point relating $\ell_0$ to $M_*$ carries no anomalous $O(1)$ coefficient beyond the derived $\sqrt2$ --- i.e.\ no large multiplicative renormalization between the bare UV stiffness and the laboratory $G$. The $0.022\%$ residual bounds any such coefficient to $1\pm2\times10^{-4}$. With that single bounded caveat, the chain $\text{shared vacuum stiffness}\to M_*\to n_0\to\ell_0\to G$ is a structural consequence of the single-vacuum architecture rather than a coincidence of two independently calibrated sectors.
\end{remark}

\section{Quantum gravity as sum over causal fillings}

Quantum gravity is not introduced by quantizing a pre-existing metric. It appears because causal fillings are not unique. For a causal boundary $\Gamma$,
\begin{equation}
\mathcal A(\Gamma)=\sum_{\Sigma:\partial\Sigma=\Gamma}\exp(iS_{\rm fill}[\Sigma]).
\end{equation}
The stationary phase approximation gives classical geometry. Non-stationary fillings give quantum gravitational corrections.
The Planck scale appears when competing fillings differ by one action quantum:
\begin{equation}
\Delta S_{\rm fill}=2\pi.
\end{equation}
Thus Planck area is interpreted as the minimal resolvable causal filling action, not as a primitive lattice unit.

\section{Matter coupling and relation to spinorial defects}

Matter enters through defects or field modes living in the causal information geometry. The stress tensor is the response of matter action to the emergent metric:
\begin{equation}
T_{\mu\nu}=-\frac{2}{\sqrt{-g}}\frac{\delta S_m}{\delta g^{\mu\nu}}.
\end{equation}
The spinorial defect sector provides a concrete candidate matter sector: spin arises from nontrivial frame holonomy, charge from residual compact $U(1)$ structure, and mass from spectral/RG closure. In the gravity programme, that matter sector supplies $T_{\mu\nu}$ while causal filling supplies the gravitational response.

The separation is important. The electron sector can be internally closed without solving gravity, while gravity can be developed as the causal-information stiffness theory into which the spinorial matter sector couples.


\section{Horizon consequences: black-hole entropy and cosmological zero mode}

The preceding sections derive the local stiffness form of gravity. Horizon physics provides two nonlocal checks: black-hole entropy probes causal filling boundaries at finite area, while the cosmological constant probes the largest causal filling boundary of the observable universe.

\subsection{Gravitational time dilation}

In this framework, clocks measure local causal update rate. Mass-energy increases causal filling stiffness. Therefore the lapse $N(x)$ decreases in high-stiffness regions:
\begin{equation}
N(x)\downarrow \quad\Longleftrightarrow\quad \kappa_{\rm caus}(x)\uparrow.
\end{equation}
Time dilation is therefore not ``time bending'' as a primitive effect; it is the slowing of local causal information transport.

\subsection{Black holes as causal filling boundaries}

A black hole is a causal filling singularity: causal transport stiffness diverges, minimal fillings become nonextendable, and information escape cost becomes horizon-localized. The event horizon is therefore treated as a causal filling boundary. Entropy counts admissible microscopic boundary fillings rather than bulk volume states.

Let $A_H$ be the horizon area and let $\ell_*^2$ be the effective causal filling area. The naive number of elementary area cells is
\begin{equation}
N_H=\frac{A_H}{\ell_*^2}.
\end{equation}
However, the physical causal filling degree of freedom is not a scalar pixel. The gravitational boundary mode is a causal bivector mode. At each horizon generator there are two null sheets and two transverse screen directions. Hence the local causal filling degeneracy is
\begin{equation}
\dim \C_p=2_{\rm null\ sheets}\times2_{\rm screen}=4.
\end{equation}
Therefore one physical horizon bit occupies four elementary causal area cells:
\begin{equation}
N_{\rm phys}=\frac{A_H}{4\ell_*^2}.
\end{equation}
The dimensionless entropy is the number of admissible physical boundary states,
\begin{equation}
\frac{S_{\rm BH}}{k_B}=N_{\rm phys}=\frac{A_H}{4\ell_*^2}.
\end{equation}
Restoring $\ell_*^2=\hbar G/c^3$ gives
\begin{equation}
\boxed{S_{\rm BH}=\frac{k_Bc^3A_H}{4\hbar G}.}
\end{equation}
This is the Bekenstein-Hawking area law, normally stated as one quarter of the horizon area in Planck-area units \cite{bekensteinhawking-scholarpedia}.

\begin{proposition}[Black-hole entropy from causal filling state count]
If black-hole horizon states are causal boundary filling states and the physical gravitational boundary mode is the four-dimensional causal bivector mode $\C_p$, then the entropy is
\begin{equation}
S_{\rm BH}=\frac{k_BA_H}{4\ell_*^2}=\frac{k_Bc^3A_H}{4\hbar G}.
\end{equation}
The factor $1/4$ is not inserted thermodynamically; it is the inverse of the causal bivector degeneracy $\dim\C_p=4$.
\end{proposition}

\subsection{Cross-verification of Newton coupling from entropy}

The black-hole entropy relation can be inverted to recover the same Newton-coupling normalization obtained from causal stiffness. From
\begin{equation}
S_{\rm BH}=\frac{k_Bc^3A_H}{4\hbar G},
\end{equation}
one obtains
\begin{equation}
\boxed{G=\frac{k_Bc^3A_H}{4\hbar S_{\rm BH}}.}
\end{equation}
Equivalently, using $S_{\rm BH}/k_B=A_H/(4\ell_*^2)$,
\begin{equation}
\ell_*^2=\frac{A_H}{4(S_{\rm BH}/k_B)},\qquad
\boxed{G=\frac{c^3}{\hbar}\ell_*^2.}
\end{equation}
Thus black-hole thermodynamics and local causal filling stiffness give the same conversion between geometric area and Newton coupling.

Combining this with the electron-sector inversion gives a non-circular cross-sector check. The electron relation
\begin{equation}
m_e=2M_*\exp\left(-\frac{32\pi^3}{19}\right)
\end{equation}
fixes
\begin{equation}
M_* = \frac{m_e}{2}\exp\left(\frac{32\pi^3}{19}\right).
\end{equation}
Using $m_ec^2=0.51099895069\,{\rm MeV}$ \cite{nist-electron-energy} gives
\begin{equation}
M_*=1.2210249\times10^{19}\,{\rm GeV}.
\end{equation}
Then
\begin{equation}
G_{\rm pred}=\frac{\hbar c^5}{M_*^2}=6.67283\times10^{-11}\,{\rm m^3\,kg^{-1}\,s^{-2}},
\end{equation}
compared with the CODATA value
\begin{equation}
G_{\rm obs}=6.67430\times10^{-11}\,{\rm m^3\,kg^{-1}\,s^{-2}}\cite{nist-g}.
\end{equation}
The relative deviation is
\begin{equation}
\frac{G_{\rm pred}-G_{\rm obs}}{G_{\rm obs}}\simeq -2.21\times10^{-4},
\end{equation}
about $0.022\%$ low. This small residual is not claimed to be explained here; plausible sources include two-loop RG terms, threshold corrections, electroweak embedding, and the distinction between bare UV stiffness and the laboratory low-energy value of $G$.

\subsection{Cosmological constant as the IR causal filling zero mode}

The derivative expansion includes a zeroth-order term
\begin{equation}
S_0=\int d^4x\sqrt{-g}\,\Lambda_{\rm eff}.
\end{equation}
In the causal filling interpretation, this term is the residual zero-mode filling curvature of the largest causal diamond. For an asymptotically de Sitter causal horizon of radius $R_C$, the curvature radius fixes
\begin{equation}
\boxed{\Lambda=\frac{3}{R_C^2}.}
\end{equation}
In flat $\Lambda$CDM the de Sitter causal radius associated with the dark-energy fraction is
\begin{equation}
R_C=\frac{c}{H_0\sqrt{\Omega_\Lambda}},
\end{equation}
so
\begin{equation}
\boxed{\Lambda=\frac{3H_0^2\Omega_\Lambda}{c^2}.}
\end{equation}
Using Planck 2018 late-Universe parameters $H_0=(67.4\pm0.5)\,{\rm km\,s^{-1}\,Mpc^{-1}}$ and $\Omega_m=0.315\pm0.007$, hence $\Omega_\Lambda\simeq0.685$ in flat $\Lambda$CDM \cite{planck2018}, gives
\begin{equation}
\Lambda_{\rm CFT}=1.0909\times10^{-52}\,{\rm m^{-2}}.
\end{equation}
The closest standard tabulated empirical value is the PDG astrophysical constant
\begin{equation}
\Lambda_{\rm PDG}=1.088(30)\times10^{-56}\,{\rm cm^{-2}}
=1.088(30)\times10^{-52}\,{\rm m^{-2}}\cite{pdg-astro2024}.
\end{equation}
The fractional difference against the central PDG value is
\begin{equation}
\frac{\Lambda_{\rm CFT}-\Lambda_{\rm PDG}}{\Lambda_{\rm PDG}}
\simeq 2.7\times10^{-3},
\end{equation}
about $0.27\%$ high and well within the quoted PDG uncertainty.

\begin{proposition}[Cosmological constant as causal IR curvature]
If the cosmic zero mode is the curvature of the largest de Sitter-like causal filling boundary, then
\begin{equation}
\Lambda=3/R_C^2.
\end{equation}
For a flat $\Lambda$CDM causal radius $R_C=c/(H_0\sqrt{\Omega_\Lambda})$, this becomes
\begin{equation}
\Lambda=3H_0^2\Omega_\Lambda/c^2,
\end{equation}
which evaluates to $1.0909\times10^{-52}\,{\rm m^{-2}}$ for Planck 2018 parameters, matching most closely the PDG value $1.088(30)\times10^{-52}\,{\rm m^{-2}}$.
\end{proposition}

\begin{remark}[Status of the cosmological constant]
This is not yet a parameter-free microscopic derivation of $\Lambda$. The local theory derives how a cosmic causal boundary sources a zeroth-order curvature term. The exact value still requires deriving $R_C$, $H_0$, or $\Omega_\Lambda$ from the event process. The result nevertheless reframes the cosmological constant as an IR causal-boundary quantity rather than a UV vacuum-energy density.
\end{remark}

\section{Predictions and falsifiable targets}

The programme becomes physically decisive only when it produces nontrivial consequences beyond recovering GR. Candidate targets are:
\begin{enumerate}[label=(\roman*)]
\item derive $G$ from a discrete causal filling spectrum or cross-sector UV scale;
\item recover the graviton propagator from fluctuations of $S_{\rm fill}$;
\item derive the Bianchi identity from causal transport conservation;
\item derive black-hole entropy from the causal filling state count and test higher-curvature corrections to the area law;
\item compute higher-curvature coefficients $c_2,c_3$ from the causal kernel;
\item predict deviations from GR in high-curvature or high-information-stiffness regimes;
\item derive Lorentzian signature uniqueness from positivity and causal irreversibility constraints;
\item derive the cosmic causal radius $R_C$ rather than using $H_0$ and $\Omega_\Lambda$ as empirical inputs.
\end{enumerate}

\section{Main formal result}

\begin{theorem}[Causal filling gravity skeleton]
Let $(\E,\pre,K,\rho)$ be a locally finite causal information structure with compact causal diamonds, positive causal Fisher kernel, stable coarse-grained filling spectrum, and a nondegenerate continuum Hessian. Then, under the stated localizability and isotropy assumptions:
\begin{enumerate}[label=(\roman*)]
\item causal order is primitive: $x\pre y$ precedes metric time;
\item time is the generator of the causal transport semigroup, $\partial_t=L$;
\item the arrow of time is the orientation of positive causal entropy production;
\item quantum phase dynamics arise from a symplectic lift $(\rho,S)\in T^*\mathcal P$;
\item paired causal update/reconstruction semigroups generate a second-order hyperbolic continuum operator;
\item Lorentzian signature follows from the split between one generator direction and spatial Fisher fluctuation directions;
\item the emergent metric is the Hessian of causal information transport;
\item curvature is infinitesimal holonomy of causal Fisher transport;
\item the causal filling action has leading curvature term $(\kappa_{\rm grav}/16\pi)\int\sqrt{-g}R$;
\item variation of causal filling plus matter gives $G_{\mu\nu}+\Lambda g_{\mu\nu}=8\pi G T_{\mu\nu}$ when $\kappa_{\rm grav}=G^{-1}$;
\item the factor $16\pi$ is derived from the minimal causal-curvature angular volume $4\pi\times4$;
\item black-hole entropy follows as $S_{\rm BH}=k_BA_H/(4\ell_*^2)=k_Bc^3A_H/(4\hbar G)$ from causal boundary state counting;
\item the cosmological constant is the IR causal-boundary curvature $\Lambda=3/R_C^2$, reducing in flat $\Lambda$CDM to $3H_0^2\Omega_\Lambda/c^2$.
\end{enumerate}
\end{theorem}

\section{Open obligations}

The programme is not complete. The remaining obligations are:
\begin{enumerate}
\item \textbf{Microscopic kernel.} Derive the causal kernel $K(x,y)$ from a discrete event process rather than assume it.
\item \textbf{Microscopic continuum theorem.} The local Fisher-Hessian continuum limit is proven under finite-moment, smooth-density, and stable-kernel assumptions. The remaining task is to derive these hypotheses from the primitive discrete event dynamics rather than impose them as coarse-graining conditions.
\item \textbf{Signature theorem.} Prove uniqueness of $(-,+,+,+)$ from causal positivity, irreversibility, and reversible fluctuation sectors.
\item \textbf{Hyperbolicity theorem.} Prove that paired causal semigroups uniquely generate the Lorentzian wave operator rather than a broader class of second-order operators.
\item \textbf{Causal Lichnerowicz spectrum.} Construct $\mathcal L_{\rm grav}$ explicitly and compute $\lambda_{\rm grav,min}$.
\item \textbf{Newton constant.} Derive $G$ from $\lambda_{\rm grav,min}/S_{\rm curv}$, not from empirical normalization.
\item \textbf{Bianchi identity.} Derive $\nabla^\mu G_{\mu\nu}=0$ from causal transport conservation.
\item \textbf{Matter coupling.} Prove that spinorial defect matter couples universally through the emergent metric.
\item \textbf{Graviton limit.} Recover the linearized spin-2 graviton propagator and gauge redundancy.
\item \textbf{Black-hole entropy corrections.} The leading area law is derived from causal boundary state counting; derive logarithmic and higher-curvature corrections from the filling spectrum.
\item \textbf{Cosmological constant.} The IR relation $\Lambda=3/R_C^2$ is derived; derive the cosmic causal radius $R_C$ from the discrete event process rather than using cosmological parameters.
\item \textbf{Higher-curvature coefficients.} Derive $c_2,c_3,\ldots$ from the causal kernel rather than fit them.
\end{enumerate}


\section{Conclusion}

Causal filling gravity proposes that gravity is the long-wavelength stiffness theory of causal information transport. The primitive data are events, causal order, a correlation kernel, and density. Time is the infinitesimal generator of coarse-grained causal updates. The arrow of time is causal entropy production. Quantum phase dynamics arise from a symplectic lift of density flow. Lorentzian hyperbolicity arises from paired forward/backward causal semigroups. The metric is the Hessian of causal information transport. Curvature is Fisher holonomy. Einstein's equation is the stationarity condition of the causal filling functional. The normalization $16\pi$ comes from the minimal causal-curvature angular volume: the celestial sphere of null directions times the four-dimensional causal bivector filling subspace.

Version 0.4 adds two horizon-scale consolidations. First, black-hole entropy is recovered from causal boundary filling states. Because the physical gravitational boundary mode is a four-dimensional causal bivector mode, one physical horizon state occupies four elementary causal filling cells. This gives
\begin{equation}
S_{\rm BH}=\frac{k_BA_H}{4\ell_*^2}=\frac{k_Bc^3A_H}{4\hbar G},
\end{equation}
matching the Bekenstein-Hawking area law and independently cross-verifying the same conversion $G=c^3\ell_*^2/\hbar$ obtained from local causal stiffness. Second, the cosmological constant is reframed as the IR zero-mode curvature of the largest causal filling boundary. The relation
\begin{equation}
\Lambda=\frac{3}{R_C^2}=\frac{3H_0^2\Omega_\Lambda}{c^2}
\end{equation}
returns $1.0909\times10^{-52}\,{\rm m^{-2}}$ for Planck 2018 parameters, closest to the PDG value $1.088(30)\times10^{-52}\,{\rm m^{-2}}$.

The framework should be understood as gravity-as-information before it is understood as quantum gravity. Quantum gravity appears as the amplitude theory of competing causal fillings. The decisive remaining tasks are now sharpened: derive the primitive causal correlation length $\ell_0$ directly from the event process, derive the cosmic causal radius $R_C$, and compute the correction hierarchy beyond the leading causal filling action. If those steps succeed, Newton's constant, the Planck scale, black-hole entropy, the cosmological constant, and higher-curvature corrections become consequences of one causal information structure rather than independent inputs.

\subsection*{Parameter economy relative to the Standard Model and general relativity}

The Standard Model carries on the order of nineteen conventional free parameters (nine charged-fermion masses, four CKM entries, three gauge couplings, two Higgs-sector parameters, the QCD vacuum angle), extended to roughly twenty-six to twenty-eight once neutrino masses and mixings are included; general relativity adds $G$ and $\Lambda$. All are fixed by measurement and none is predicted from within the theory.

In the present framework, by contrast, all gauge couplings, the Weinberg angle, the ultraviolet fine-structure constant, the generation count (given $\alpha_{\mathrm{em}}(m_e)$), the electron-to-Planck mass ratio, and Newton's constant are derived from group theory and spectral geometry with no adjustable parameter. At the present stage the framework retains: (i) one absolute mass scale ($\ell_0$, equivalently $M_*$, $m_e$, or $G$), which sets units and is a candidate for derivation from the exponent-vector kernel of the Arithmetic Vacuum (\emph{Physical Sector}, Open Problem 11.2); (ii) $\alpha_{\mathrm{em}}(m_e)$, used to fix the generation count, a candidate for derivation from the topological classification of stable spinorial defect sectors (\emph{Fisher-Frame}, Open Problem 1); (iii) the observed $(3{+}1)$-dimensional spacetime, used to select the coarse-graining sector $k=2$, a candidate for derivation by the internalisation principle (\emph{Physical Sector}, Open Problem 11.4); and (iv) the electroweak symmetry-breaking scale (the Higgs VEV $v\approx246$~GeV) together with the cosmological parameters $H_0,\Omega_\Lambda$ entering the IR sector, which currently require separate dynamical input. The first three are structurally connected to the Arithmetic Vacuum and admit candidate derivations by upstream work; only the electroweak-breaking scale requires genuinely separate dynamics. Where the Standard Model and general relativity together fix of order twenty-six dimensionless parameters by hand, the present framework fixes one unit-setting scale and at most two further empirical inputs, both of which are candidates for elimination.

\begin{thebibliography}{9}

\bibitem{nist-electron-energy}
NIST/CODATA, ``Electron mass energy equivalent in MeV,'' The NIST Reference on Constants, Units, and Uncertainty. Numerical value $m_ec^2=0.51099895069\,{\rm MeV}$. Available at \url{https://physics.nist.gov/cgi-bin/cuu/Value?mec2mev=}.

\bibitem{nist-hbar}
NIST/CODATA, ``Reduced Planck constant,'' The NIST Reference on Constants, Units, and Uncertainty. Numerical value $\hbar=1.054571817\times10^{-34}\,{\rm J\,s}$, exact in the present SI. Available at \url{https://physics.nist.gov/cgi-bin/cuu/Value?hbar=}.

\bibitem{nist-g}
NIST/CODATA, ``Newtonian constant of gravitation,'' The NIST Reference on Constants, Units, and Uncertainty. Numerical value $G=6.67430\times10^{-11}\,{\rm m^3\,kg^{-1}\,s^{-2}}$. Available at \url{https://physics.nist.gov/cgi-bin/cuu/Value?bg=}.


\bibitem{bekensteinhawking-scholarpedia}
Scholarpedia, ``Bekenstein-Hawking entropy.'' Summarizes the entropy as one quarter of the horizon area in Planck-area units. Available at \url{https://www.scholarpedia.org/article/Bekenstein-Hawking_entropy}.

\bibitem{planck2018}
Planck Collaboration, ``Planck 2018 results. VI. Cosmological parameters,'' \emph{Astronomy \& Astrophysics} 641, A6 (2020). Reports $H_0=(67.4\pm0.5)\,{\rm km\,s^{-1}\,Mpc^{-1}}$ and $\Omega_m=0.315\pm0.007$ in base $\Lambda$CDM. Available at \url{https://arxiv.org/abs/1807.06209}.

\bibitem{pdg-astro2024}
Particle Data Group, ``Astrophysical Constants and Parameters,'' Review of Particle Physics 2024/2025. Reports $\Omega_\Lambda=0.685(7)$ and $\Lambda=1.088(30)\times10^{-56}\,{\rm cm^{-2}}$. Available at \url{https://pdg.lbl.gov/2025/reviews/rpp2024-rev-astrophysical-constants.pdf}.

\end{thebibliography}

\appendix
\section{Symbol glossary}

\begin{center}
\begin{tabular}{ll}
\toprule
Symbol & Meaning\\
\midrule
$\E$ & Set of elementary events\\
$\pre$ & Causal preorder\\
$K(x,y)$ & Causal correlation kernel\\
$\rho(x)$ & Event density\\
$\Gamma$ & Causal cycle in the order complex\\
$\Sigma$ & Causal filling with $\partial\Sigma=\Gamma$\\
$S_{\rm fill}$ & Causal filling action/cost\\
$\kappa_{\rm caus}$ & Causal filling modulus\\
$\tau_\Gamma$ & Intrinsic causal diameter\\
$S_{\rm info}$ & Causal information action\\
$g_{\mu\nu}$ & Emergent metric Hessian\\
$T_\tau$ & Causal transport semigroup\\
$L$ & Infinitesimal causal generator\\
$\partial_t$ & Emergent time derivative $L$\\
$\Sigma[\rho,K]$ & Causal entropy production\\
$D_+,D_-$ & Forward/backward Nelson-type derivatives\\
$\Box_g$ & Emergent Lorentzian wave operator\\
$\Omega$ & Fisher transport connection\\
$U_\Gamma$ & Causal holonomy around $\Gamma$\\
$R_{\mu\nu\rho\sigma}$ & Emergent Riemann curvature\\
$\C_p$ & Causal filling bivector subspace\\
$\Omega_{\rm grav}$ & Minimal causal-curvature angular volume $16\pi$\\
$\kappa_{\rm grav}$ & Gravitational filling stiffness $G^{-1}$\\
$\ell_0$ & Primitive microscopic causal correlation length\\
$n_0$ & Primitive event/correlation-cell density; conditional closure $n_0=4M_*^4$\\
$M_*$ & Shared UV vacuum stiffness scale probed by charged spinorial electron sector\\
$C_F$ & Dimensionless causal kernel second moment, $C_F=1/2$ at the Gaussian fixed point\\
$\ell_P$ & Effective Planck/filling length, $\ell_P=\sqrt2\,\ell_0$\\
$\mathcal L_{\rm grav}$ & Causal Lichnerowicz/filling operator\\
$\lambda_{\rm grav,min}$ & Lowest nontrivial gravitational filling eigenvalue\\
$G$ & Newton coupling, inverse gravitational stiffness\\
$N,N^i$ & Lapse and shift in the continuum ADM limit\\
\bottomrule
\end{tabular}
\end{center}

\section{Derived versus assumed ingredients}

\begin{center}
\small
\renewcommand{\arraystretch}{1.15}
\begin{tabularx}{\textwidth}{>{\raggedright\arraybackslash}p{0.22\textwidth}>{\raggedright\arraybackslash}p{0.30\textwidth}X}
\toprule
Item & Status & Comment\\
\midrule
Causal order primitive & Assumption & Starting ontology\\
Metric as Hessian & Derived under continuum nondegeneracy & Needs rigorous limit\\
Time as generator & Derived from semigroup consistency & $T_\tau=e^{\tau L}$\\
Arrow of time & Derived from entropy production & Requires asymmetric kernel/coarse-graining\\
Quantum phase & Derived from symplectic lift & Lift is structural assumption\\
Lorentzian operator & Derived from paired semigroups plus cone constraint & Needs uniqueness proof\\
Einstein-Hilbert form & Derived from local derivative expansion & Standard uniqueness at two derivatives\\
Einstein equation & Derived by variation & Given coefficient normalization\\
$16\pi$ & Derived from $4\pi\times4$ & Needs formalization of causal bivector trace\\
$G$ coefficient relation & Derived under Gaussian fixed point & $G=2\ell_0^2$ in natural units\\
$G$ absolute SI value & Conditional cross-sector result & Electron-sector inversion predicts $G$ within $0.022\%$; v0.7 closes $\ell_0$ if the charged spinorial electron sector probes the shared UV vacuum stiffness\\
Event density $n_0$ & Conditional microscopic closure & $n_0=4M_*^4$ follows from Poisson large deviations plus cross-sector stiffness matching; primitive-event identification remains a hypothesis\\
Black-hole entropy & Leading law derived & $S_{\rm BH}=k_BA_H/(4\ell_*^2)$ from causal boundary state count; corrections open\\
Cosmological constant & IR relation derived & $\Lambda=3/R_C^2$; deriving $R_C$ from event microphysics remains open\\
\bottomrule
\end{tabularx}
\end{center}

\end{document}
